\documentclass[11pt]{article}
\usepackage[a4paper,margin=1in]{geometry}
\usepackage{amsmath,amssymb,amsthm}
\usepackage{booktabs}
\usepackage{array}
\usepackage{hyperref}

\newtheorem{definition}{Definition}
\newtheorem{theorem}{Theorem}
\newtheorem{conjecture}{Conjecture}

\newcommand{\CRCFT}{\mathrm{CRCFT}}
\newcommand{\CTMT}{\mathrm{CTMT}}
\newcommand{\CRE}{\mathrm{CRE}}
\newcommand{\XiC}{\Xi_C}

\title{Step 183: Classical-RH Carrier Foreclosure Taxonomy}
\author{RH Membrane Adequacy Track}
\date{}

\begin{document}
\maketitle

\section*{Verdict}

\[
\boxed{\mathrm{final\_verdict}=V_{\mathrm{CRCFT\_formalized}}.}
\]

This step formalizes the Classical-RH Carrier Foreclosure Taxonomy
(\(\CRCFT\)) as a candidate foundational typed condition.  It is not a branch
closeout and not an RH proof.  It is a taxonomy for the typed ways in which
classical RH-equivalent carriers terminate under the framework discipline.

\begin{definition}[Classical-RH-equivalent carrier]
A typed primary carrier \(C\) carrying a residual \(\XiC\) aimed at RH is
classically-RH-equivalent, abbreviated \(\CRE\), when either:
\begin{enumerate}
\item \(C\) is one of the named classical RH approaches, such as
Hilbert--Polya, de Branges, Connes adelic/semilocal, Selberg trace,
Burnol/Sonine, or Hecke carriers; or
\item closure of \(\XiC\) is logically equivalent to RH.
\end{enumerate}
\end{definition}

\begin{definition}[Classical-RH Carrier Foreclosure Taxonomy]
A \(\CRE\) carrier exhibits \(\CRCFT\) when, after the framework's legality,
carrier, and no-smuggling audits are applied, the attempted closure terminates
in at least one of the following typed modes:
\[
\CRCFT\text{-}\CTMT,\qquad
\CRCFT\text{-TE},\qquad
\CRCFT\text{-BF}.
\]
\end{definition}

\section*{Modes}

\begin{center}
\begin{tabular}{p{0.18\linewidth}p{0.24\linewidth}p{0.20\linewidth}p{0.28\linewidth}}
\toprule
Mode & Name & Obstruction level & Framework consequence\\
\midrule
\(\CRCFT\)-\(\CTMT\) & matrix-element terminus
& post-reduction
& identify the carrier-native matrix-element family and its resolution mode\\
\(\CRCFT\)-TE & target-equivalence
& pre-reduction
& central positivity, spectral, or structural condition is RH-equivalent\\
\(\CRCFT\)-BF & bridge-failure
& carrier-comparison
& carrier is non-comparable to the target ledger or has wrong spectral type\\
\bottomrule
\end{tabular}
\end{center}

\section*{Instances}

\begin{center}
\begin{tabular}{p{0.22\linewidth}p{0.20\linewidth}p{0.19\linewidth}p{0.27\linewidth}p{0.07\linewidth}}
\toprule
Carrier & Route & Mode & Specific obstruction & Step\\
\midrule
Burnol/Sonine & Branch A Calkin G2--G5
& \(\CRCFT\)-\(\CTMT\)-stuck
& \(\kappa\) vector / Burnol projected-kernel resolvent theorem
& 179\\
Burnol/Sonine & Branch B SL164.1
& \(\CRCFT\)-\(\CTMT\)-stuck
& same \(\kappa\) vector / projected Sonine kernel theorem
& 177--178\\
Burnol/Sonine & Branch C shifted co-Poisson
& \(\CRCFT\)-\(\CTMT\)-foreclosed-numerical
& nonzero \(L_{\rho,k}(G)\) across four data points, \(|L|\ge 0.11\)
& 175--176\\
Hecke carrier & H6 zeta-fiber descent
& \(\CRCFT\)-BF
& V-NC ledger-relative non-comparability
& 168\\
de Branges \(H(E_{\rm RH})\)
& \(\Xi_{\rm dB}\) closure
& \(\CRCFT\)-TE
& de Branges RH condition and Conrey--Li survival
& 181\\
HP/BK standard \(H_{xp}\)
& standard spectral ledger
& \(\CRCFT\)-BF
& continuous \(\mathbb R\) spectrum; interval spectra arithmetic
& 182\\
HP/BK modified
& cutoff/boundary exact matching
& \(\CRCFT\)-TE
& exact matching is a non-inherited HP/RH assertion
& 182\\
\bottomrule
\end{tabular}
\end{center}

\begin{theorem}[CRCFT consequence]
When a \(\CRE\) carrier \(C\) exhibits \(\CRCFT\) in mode
\[
m\in\{\CTMT,\mathrm{TE},\mathrm{BF}\},
\]
the framework output is the precise typed obstruction associated with that
mode:
\begin{itemize}
\item for \(m=\CTMT\), the carrier-native matrix-element family and its
resolution mode;
\item for \(m=\mathrm{TE}\), the target-equivalence statement showing that
direct pursuit is structurally circular;
\item for \(m=\mathrm{BF}\), the bridge-failure statement naming the failed
carrier-comparison axis.
\end{itemize}
\end{theorem}

\begin{proof}
This is a synthesis theorem over the recorded instances.  In the Burnol/Sonine
case the operator reductions reach explicit matrix-element families
\(\kappa\) or \(L_{\rho,k}(G)\); these are \(\CTMT\) outputs.  In the
de Branges and modified HP/BK cases, the central condition is already an
RH-equivalent structural or spectral assertion, so the framework output is the
target-equivalence no-go.  In the Hecke H6 and standard HP/BK cases, the
carrier comparison itself fails: either the ledgers are non-comparable or the
operator has the wrong spectral type.  These three cases exhaust the
obstruction classes represented by the seven instances and give the stated
typed outputs.
\end{proof}

\begin{conjecture}[CRCFT coverage]
Every classically-RH-equivalent carrier exhibits \(\CRCFT\) in at least one of
the three modes \(\CTMT\), TE, or BF.  No \(\CRE\) carrier is
closure-amenable without supplying external content matching at least one of:
a \(\CTMT\)-resolving classical theorem, a \(\CTMT\)-foreclosing
counterexample, an RH resolution itself for TE, or a non-inherited carrier
bridge for BF.
\end{conjecture}

This is a conjecture, not a theorem.  Its current support is 4 carriers and 7
instances.  It is refutable by exhibiting a \(\CRE\) carrier with a completed,
non-target-equivalent, non-\(\CRCFT\) closure.

\section*{Relationship to CTMT}

\(\CTMT\) is the matrix-element-terminus typed condition formalized at Step
180.  \(\CRCFT\) is broader: it contains \(\CTMT\) as the post-reduction mode
and adds target-equivalence and bridge-failure for pre-reduction and
carrier-comparison failures.  The two typed conditions are therefore related
but distinct.

\section*{Corpus Candidate}

\(\CRCFT\) is flagged for future foundational-corpus inclusion, either as a
sibling file \texttt{anti\_loc/CRCFT.md} or as an extension to the Step 180
\texttt{CTMT} note.  Integration is not performed in this step.

\section*{Nonclaim Boundary}

\(\CRCFT\) does not prove RH, does not close any residual, and does not remove
any retained no-go.  The coverage statement is conjectural.  The taxonomy
names typed obstruction interfaces; it is not itself a closure mechanism.

\end{document}
